\section{Decoder-Time Hard Masking for Trusted-Model Privacy}
\label{sec:hard-mask}

We implement a decoder-time hard-masking path for the trusted local model.  The
setting is a two-party agent interaction: an external or attacker model provides
an adversarial user message, while the trusted model receives both this message
and a private source document.  The source document is placed in the trusted
model's system context, whereas the attacker text is placed in the user message.
Protected attributes are supplied explicitly as a list of values or structured
records, rather than being inferred from the benchmark labels.  This separation
is important for future POLAR-style evaluation: the defense should operate on a
policy supplied to the trusted runtime, while gold protected values are reserved
for debugging and scoring.

\paragraph{Local trusted-model path.}
The current implementation uses a local Hugging Face model, with
\texttt{swiss-ai/Apertus-8B-Instruct-2509} as the default model identifier and
\texttt{models/apertus-8b-instruct-2509} as the default local path.  Local
execution is necessary because true decoder-time hard masking requires access to
the next-token logits.  Hosted chat APIs that do not expose logits can support
post-hoc checking or rewriting, but they cannot implement the same hard
generation-time constraint.

\paragraph{Exact token-level masking.}
The first masking mechanism follows the constrained-decoding view of privacy
defense: before a token is selected, the runtime checks whether appending each
candidate vocabulary token would complete a protected string or match a private
regex pattern, and sets the corresponding logit to $-\infty$ if so.  Let
$y_{<t}$ be the generated assistant prefix, $v$ a candidate token, and
$\tau(v)$ its decoded text.  For each protected surface form
$\gamma \in \Gamma(A_{\mathrm{protected}})$, the runtime blocks $v$ whenever
$y_{<t}\tau(v)$ would complete $\gamma$.  This is implemented by
\texttt{HFPrivacyLogitsProcessor}, which only inspects tokens generated after
the prompt boundary, so protected values that appear in the private source
document do not count as already generated assistant output.

This mechanism is closest to constrained decoding for privacy-preserving LLM
inference, where logits are masked to prevent structured PII from being emitted
\cite{constrained_decoding_privacy_llm}.  Our implementation differs in that the
blocked strings are policy-supplied protected attributes and surface forms, not
only fixed regex categories such as emails or phone numbers.  The same module
can still include regex constraints for common private formats.

\paragraph{Why prefix blocking was not enough.}
Early experiments showed that aggressive prefix blocking can over-mask benign
tokens.  For example, if the protected value is \texttt{\$2M}, shell quoting
mistakes may reduce it to \texttt{M}; a prefix-aware mask then blocks many
ordinary tokens containing \texttt{m}.  Even when the value is quoted correctly,
blocking partial prefixes such as \texttt{alice@example.} can make generation
unnatural and can force the model to terminate early.  For this reason, the
current demo does not rely on partial-prefix blocking.  It allows the model to
temporarily generate a protected value internally, detects the full leakage in
the assistant text, rewinds the live generation state, and regenerates from an
earlier state.

\paragraph{Rewind-on-leak masking.}
The main hard-mask mode is \emph{rewind-on-leak}.  Generation proceeds greedily
one token at a time.  After each selected token, the runtime decodes only the
assistant-generated tokens and checks whether any full protected value appears
case-insensitively.  If no protected value is present, generation continues.  If
a full protected value appears, the runtime finds a rewind point, removes all
generated tokens after that point, records the token that opened the leaking
branch, bans that token at the rewind state, and continues decoding.  State bans
are stored as
\[
    \mathcal{B}: \mathrm{tuple}(y_{1:k}) \mapsto \{v_1,\ldots,v_m\},
\]
so a token is banned only at the specific generation state that previously led
to leakage.  This keeps the intervention local: a token such as
\texttt{Alice}, \texttt{al}, or a comma may be banned at one state while
remaining available elsewhere.

The default rewind strategy is \texttt{value}.  If a protected value starts at
character position $i$ in the decoded assistant text, the runtime estimates the
first token whose decoded prefix crosses $i$, rewinds to the token before that
point, and bans the leaking token at that state.  This strategy is simple and
strong for exact-value leakage, but it can be too late when the model has
already opened a syntactic slot such as \texttt{His email is \_} or
\texttt{The patient, \_}.  In those cases, regenerating from immediately before
the value often causes the model to refill the same private slot.

\paragraph{Dependency-role rewind.}
To address slot-refilling failures, we add an optional
\texttt{dependency} rewind strategy.  The dependency parser is invoked only
after a complete protected value has already appeared in the generated assistant
text, so it does not run at every token step.  If spaCy or the configured model
is unavailable, the runtime falls back to rule-based line, clause, or value
rewind and records the fallback in the trace.

The dependency strategy chooses an earlier rewind point according to the role of
the protected span:
\begin{itemize}
    \item \textbf{Field or markdown line.}  If the protected value appears in a
    line such as \texttt{Email: x}, \texttt{Phone: x}, or
    \texttt{Patient: x}, the runtime rewinds to the beginning of that line.
    This rule is applied before dependency parsing because parser analyses of
    markdown lists are often unstable.
    \item \textbf{Subject-like value.}  If the protected value is a subject
    (\texttt{nsubj}, \texttt{nsubjpass}, or \texttt{csubj}), the runtime rewinds
    to the current sentence or clause start.
    \item \textbf{Copular slot value.}  If the value is an attribute or
    complement in a copular construction, such as \texttt{His email is
    121@gmail.com}, the runtime rewinds to the slot subject start rather than to
    the email itself.  This removes the whole slot expression
    \texttt{His email is ...}, making it less likely that the model immediately
    fills the same private slot again.
    \item \textbf{Object or complement.}  If the value is an object or clausal
    complement, the runtime rewinds to the governing predicate.  For example,
    \texttt{The note includes H-133326} rewinds to \texttt{includes}.
    \item \textbf{Prepositional object.}  If the protected value is the object
    of a preposition, such as \texttt{Contact him at +41...}, the runtime
    rewinds to the governing predicate rather than only to the preposition.
    \item \textbf{Apposition.}  If the model generates an appositive disclosure,
    such as \texttt{The patient, Lina Roth}, the runtime rewinds to the
    punctuation that opened the apposition.  The next regeneration then bans the
    comma at the state \texttt{The patient}, discouraging the model from
    re-opening the same name slot.
    \item \textbf{Unknown role.}  If no reliable role is found, the runtime
    falls back to a clause start and then to the protected value start.
\end{itemize}

This dependency-role design is not a general semantic privacy guarantee.  It is
a decoder-time repair heuristic that tries to rewind to the structure that made
the protected value likely, rather than merely deleting the value itself.  It is
most useful for common disclosure templates in privacy benchmarks: field lists,
copular slots, appositions, and contact-information predicates.

\paragraph{Traceability and failure modes.}
The demo emits both the unmasked baseline and the hard-masked response.  When
tracing is enabled, detailed events are written to a separate JSON file.  Each
rewind event records the matched protected value, the text before rewind, the
rewind target text, the token banned at that state, and dependency metadata such
as the dependency role, governing head, parser availability, fallback strategy,
and removed text.  This makes failures inspectable.  For example, if a protected
name is repeatedly produced as \texttt{- Patient's name: ...},
\texttt{* Patient's name: ...}, and \texttt{1. Patient's name: ...}, the trace
reveals that the model is changing list markers rather than abandoning the
privacy-disclosing field.  The runtime can then be improved with structural
bans, list-context priority, or repeated-leak escalation.

The current method has three important limitations.  First, it prevents exact
protected values from appearing in the final decoded output, but it does not by
itself prevent semantic leakage, such as paraphrases or quasi-identifiers.
Second, if the model repeatedly attempts to disclose the same value through
different formats, the decoder may consume many rewind attempts before producing
a useful answer.  Third, dependency parsing can choose an overly broad rewind
point in long sentences, for example mapping a leaked list item back to a remote
verb such as \texttt{share} instead of the local phrase \texttt{which includes}.
For these cases, the next improvements are list/colon boundary priority,
maximum rewind-distance caps, repeated-leak escalation, and attribute-type-aware
slot rules.

\paragraph{Relation to prior work.}
The hard-mask component is most directly related to constrained decoding for
privacy-preserving inference, which argues for inference-time prevention rather
than post-hoc filtering once PII has already been generated
\cite{constrained_decoding_privacy_llm}.  It is also related to generation-time
secret-leakage defenses such as PRISM, which monitor and intervene during
generation in multi-agent pipelines \cite{prism_secret_leakage}.  Our runtime
adopts the same broad principle that leakage should be stopped during decoding,
but focuses on policy-supplied protected attributes and trusted-context
disclosure rather than only credentials or structured secrets.

At the system level, the design is complementary to contextual-integrity and
agent privacy systems such as PrivacyChecker, AirGapAgent, and PlanTwin
\cite{privacychecker,airgapagent,plantwin}.  Those systems emphasize deciding
what information flows are appropriate, minimizing access to private data, or
using local abstractions and budgets for cloud-assisted planning.  Our hard mask
is a lower-level enforcement mechanism: once a trusted model has private
context, the decoder prevents explicit protected values from crossing the
communication boundary.  Benchmarks such as ConfAIde and POLAR motivate this
setting because they show that privacy depends on context, task utility, and
adversarial probing rather than on static PII detection alone
\cite{confaide,polar_bench}.

% Suggested BibTeX keys used above:
% - constrained_decoding_privacy_llm:
%   Constrained Decoding for Privacy-Preserving LLM Inference,
%   OpenReview riu8VN6Do8.
% - prism_secret_leakage:
%   PRISM, generation-time mitigation of secret leakage in multi-agent LLM
%   pipelines, arXiv:2605.10614.
% - privacychecker:
%   PrivacyChecker / Privacy in Action, contextual-integrity checks for
%   LLM-powered agents, arXiv:2509.17488.
% - airgapagent:
%   AirGapAgent, privacy-conscious conversational agents and contextual
%   integrity, arXiv:2405.05175.
% - plantwin:
%   PlanTwin, privacy-preserving cloud-assisted LLM agents, arXiv:2603.18377.
% - confaide:
%   ConfAIde, contextual-integrity benchmark for LLM privacy,
%   arXiv:2310.17884.
% - polar_bench:
%   POLAR-Bench, diagnostic benchmark for privacy-utility trade-offs in LLM
%   agents.
